% Sec. 4, Discussion and Conclusion (track sim, 2026-11-13; rewritten after critic round 1,
% jpw-review/sim/critiques/disc_critique_round1.md, adjudication in jpw-review/sim/disc_round1_adjudication.md).
% Structure (LDR): a lead that states the answer and the threads; what qualifies the answer (limits as links to
% Secs. 2.3 and 3.3 and three further points; LI4); what a pseudo-threshold compares; how the schedule and the
% layout act; consequences for an experiment (hardware appears only as an outlook, CONTRACT Sec. 6; no device is
% named); a resolution that synopsizes in words, answers the question of Sec. 1 and returns to its gap (RS8).
% No numeral of Sec. 3 is repeated (CONTRACT Sec. 5.7); the statements in words are checked against the data by
% jpw-review/sim/checks_appendix.py (ids S4-*). Outline and sentence table: jpw-review/sim/discussion-outline.md.
\section{Discussion and Conclusion}\label{sec:discussion}

The simulations answer the question of \autoref{sec:intro} with a pseudo-threshold that the postselection schedule, the layout and the local dimension each move. This section states what qualifies that answer, what sets the pseudo-threshold and what follows for an experiment.

The answer holds among the accepted shots (\autoref{sub:lrb_d}), within the limits of the fits near the pseudo-thresholds (\autoref{sub:pseudo_thresholds}), and without a mechanism for the larger qubit logical error rate (\autoref{sub:qubit_comparison}). Two further points qualify it. First, a pseudo-threshold compares error rates per unit of depth, so to leading order it leaves out the loss (\autoref{sub:interval_results}). Second, because the simulations give the gates on the long-range coupler the noise of any other two-qutrit gate (\autoref{sub:ancilla_layouts}), the values of the one-ancilla layout assume a long-range coupler as good as the other couplers of the processor.

% Interpretive statements 4 to 7 and 20 to 22 of jpw-review/outline/4-results-b.md, list D (what a pseudo-threshold compares).
A pseudo-threshold characterizes a pair, the code and its baseline. Like the SD6 values of the one-ancilla layout (\autoref{sub:pseudo_thresholds}) and the qubit values (\autoref{sub:qubit_comparison}), every pseudo-threshold therefore moves with the baseline. Provided that, between the crossings with the two baselines, the logical error rate grows at least as $p^2$ and the baseline error rate at most as $p$, a baseline that compiles a Clifford into fewer primitive gates than ours (\autoref{sub:fit_model}) would lower the pseudo-threshold by no more than the factor by which the baseline error rate falls.

% Interpretive statements 1 to 3, 12, 15 of list D (how the schedule and the layout act).
The results of \autoref{sub:terminal_results} to \autoref{sub:layout_comparison} are consistent with one mechanism behind the dependence of the logical error rate, and through it of the pseudo-threshold, on the schedule and on the layout: to leading order at small $p$, pairs of faults that no postselected round separates. In this picture a longer check round adds faults to every stretch between postselected rounds, and without check rounds the stretch is the whole circuit (\autoref{sub:terminal_results}). The persisting errors of \autoref{sec:intro}~\cite{Ziyad2025} are, in these circuits, faults that no postselected round has yet separated from the next, so the departures from a single exponential, which tend to grow as fewer rounds enter the postselection (\autoref{tab:chi_largek}), fit this mechanism as well.

% Outlook for an experiment (CONTRACT Sec. 6). No number of Sec. 3 is repeated; the directions are checked (S4-out-*).
For an experiment on a qutrit processor the results have three consequences. First, near the pseudo-thresholds almost every shot is rejected beyond the first few depths (\appref{app:campaign-search}), so an experiment that locates a pseudo-threshold can fit only those depths. Second, the two noise models, which weight one-qutrit gates, idling, reset and measurement differently against the two-qutrit gates (\autoref{sub:noise_models}), give pseudo-thresholds that differ by up to somewhat more than a third (\autoref{sub:pseudo_thresholds}), so a prediction for a given processor needs its own noise model. Third, a long-range coupler that is worse than the other couplers lowers only the values of the one-ancilla layout, so on such a processor the layout gaps of \autoref{sub:layout_comparison} bound from above what the two-ancilla layout costs.

In conclusion, with postselection on the check rounds in place of error correction, a logical and a physical qutrit gate can be compared already at distance two. With every check round postselected, the logical error rate per unit of depth of the $[[5,1,2]]_3$ code grows about as $p^2$ at small $p$ and stays below the linearly growing baseline error rate up to a pseudo-threshold. This pseudo-threshold falls by more than half when only every fourth round is postselected, and without check rounds the accepted signal gives no error rate and no pseudo-threshold. Whether every round or only every second, third or fourth one is postselected, the two-ancilla layout, with its longer check round, lowers the pseudo-threshold by about a tenth to a third, and the qubit code has a pseudo-threshold lower by a factor of about two to three and a half. Among the accepted shots, and with every check round postselected, the simulations place the pseudo-threshold of the logical Clifford gates of this distance-two qutrit code, under both circuit-level noise models, at two-qutrit gate error probabilities of a few percent.
